\documentclass[11pt,a4paper]{article}

\usepackage[margin=1in]{geometry}
\usepackage{amsmath,amssymb,amsthm}
\usepackage{array}
\usepackage[hidelinks]{hyperref}

\setlength{\emergencystretch}{1.5em}

\newtheorem{theorem}{Theorem}[section]
\newtheorem{proposition}[theorem]{Proposition}
\newtheorem{lemma}[theorem]{Lemma}
\newtheorem{corollary}[theorem]{Corollary}
\theoremstyle{definition}
\newtheorem{definition}[theorem]{Definition}
\newtheorem{example}[theorem]{Example}
\theoremstyle{remark}
\newtheorem{remark}[theorem]{Remark}

\title{\textbf{Disproof and Exact Sharp Classification for Conjecture 00000004507:\\
Closure System Representations of Galois Inverses}}
\author{Formal Mathematical Resolution \& Lean 4 Verification}
\date{}

\begin{document}
\maketitle

\begin{abstract}
We resolve Conjecture \texttt{00000004507}, which asserts that for a monotone Galois connection between posets, the minimal representation of its inverse is a closure system generated by units, the representation length is an explicit count of units, and finite lattices attain the lower bound $2$. We prove that the conjecture is \textbf{false as stated} under every standard order-theoretic interpretation of Galois inverses, closure systems, unit generators, and representation length. Specifically: (1) on general posets, the image of a right adjoint need not admit meets at all (refuted on the $2$-element antichain $A_2$), and even on complete distributive lattices such as the open rounded Dedekind cuts $\mathrm{Cut}(\mathbb{D}) \cong ([0,1], \le)$, a closure system need not admit \emph{any} minimal meet-generating set; (2) on finite lattices, although every closure system $\mathcal{C} \subseteq L$ admits a unique minimum meet-generating set $\mathcal{M}(\mathcal{C})$ of meet-irreducible elements and a minimum concrete set-closure ground set of size $|\mathcal{J}(\mathcal{C})|$, the claimed lower bound $2$ fails on the $1$-element lattice $C_1$ (lengths $0$ and $1$), fails on the nontrivial $2$-element chain $C_2$ even for the bijective identity Galois connection $(\mathrm{id}_{C_2}, \mathrm{id}_{C_2})$ (length $1 < 2$), fails on every finite lattice via constant Galois connections (lengths $0$ and $1$), and fails on arbitrarily large finite lattices via non-constant $2$-closed-element Galois connections (length $1 < 2$). Finally, we establish the exact salvaged theorems and sharp thresholds: $|\mathcal{M}(\mathcal{C})| \ge 2 \iff |\mathcal{J}(\mathcal{C})| \ge 2 \iff |\mathcal{C}| \ge 3$, with equality $2$ attained up to isomorphism exclusively at $C_3$ and $B_2$. All main theorems, counterexamples, and positive witnesses are formally verified in self-contained Core Lean~4.19.0 and independently checked by exhaustive enumeration in Python.
\end{abstract}

\section{Introduction and Problem Statement}

Conjecture \texttt{00000004507} from \emph{The Last Math Competition} reads verbatim in English as follows:
\begin{quote}
\textbf{Definition:} A Galois connection is a monotone adjoint pair between posets.\\
\textbf{Conjecture:} The minimal representation of its inverse is a closure system generated by units, the representation length is an explicit count of units, and finite lattices attain the lower bound $2$. \emph{(closure representation lower bound of Galois inverses)}
\end{quote}
The parallel Chinese formulation defines a Galois connection as a monotone adjoint pair between posets (\emph{pianxuji jian de dandiao bansui dui}) and asserts that the minimal representation of its inverse is a closure system generated by units (\emph{danyuan shengcheng de bibao xitong}), the representation length is an explicit count of units (\emph{danyuan shu de xianshi jishu}), and finite lattices attain the lower bound $2$ (\emph{youxian ge dadao xiajie $2$}).

\subsection{Summary of Main Verdict and Obstructions}

The conjecture is \textbf{false as stated}. A rigorous audit reveals three distinct structural obstructions, while simultaneously uncovering the exact restricted regime in which a sharp lower bound of $2$ genuinely holds:

\begin{enumerate}
\item \textbf{Failure of the lower bound $2$ on finite lattices across all candidate measures.} Let $(f, g) : L \leftrightarrows Q$ be a monotone Galois connection out of a finite lattice $L$, with closed-element closure system $\mathcal{C} = \mathrm{Im}(g) = \mathrm{Fix}(g \circ f) \subseteq L$. Across all four standard definitions of closure-system representation length---(I) minimum meet-generating set cardinality $|\mathcal{M}(\mathcal{C})|$ under arbitrary meets (including $\bigwedge \emptyset = \top_L$), (II) minimum ground-set cardinality $|U|$ in a concrete set-closure embedding $\mathcal{C} \hookrightarrow (\mathcal{P}(U), \subseteq)$, (III) minimum nonempty-meet generator cardinality $|\mathcal{M}(\mathcal{C}) \cup \{\top_L\}| = |\mathcal{M}(\mathcal{C})| + 1$ or closed-element count $|\mathcal{C}|$, and (IV) minimum unit-implicational basis size $|\Sigma_{\min}(\mathcal{C})|$---values strictly below $2$ occur on finite lattices:
\begin{itemize}
\item On $C_1 = \{0\}$ with $(\mathrm{id}_{C_1}, \mathrm{id}_{C_1})$, Readings I, II, IV give $0 < 2$ and Reading III gives $1 < 2$.
\item Even on the \emph{nontrivial} $2$-element finite lattice $C_2 = \{0 < 1\}$ with the \emph{bijective} Galois connection $(\mathrm{id}_{C_2}, \mathrm{id}_{C_2})$, Readings I and II give $1 < 2$ and Reading IV gives $0 < 2$.
\item On \emph{every} pair of finite lattices $L_1, L_2$ (even with $|L_1|, |L_2| \ge 2$), the constant monotone Galois connection $f(x) = \bot_{L_2}$, $g(y) = \top_{L_1}$ yields $\mathcal{C} = \{\top_{L_1}\}$, giving $0 < 2$ under Readings I, II and $1 < 2$ under Reading III.
\item On \emph{every} finite lattice $L$ with $|L| \ge 2$ and any $a \neq \top_L$, there exists a \emph{non-constant} monotone Galois connection $L \leftrightarrows C_2$ with $\mathcal{C} = \{a, \top_L\}$ of Reading~I and~II length $1 < 2$.
\end{itemize}

\item \textbf{Failure of minimal unit-generated closure representations on general posets and infinite complete lattices.} Although the conjecture defines Galois connections between arbitrary posets, on non-lattice posets such as the $2$-element antichain $A_2$, $\mathrm{Im}(\mathrm{id}_{A_2})$ possesses neither an empty meet nor binary meets. More deeply, even on a complete distributive lattice---namely the lattice $\mathrm{Cut}(\mathbb{D})$ of open rounded upper Dedekind cuts of dyadic rationals in $(0,1)$---for $(\mathrm{id}, \mathrm{id})$ there are \emph{no} meet-irreducible elements and \emph{no} minimal meet-generating sets: every meet-generating family $B$ stays meet-generating after deleting an element $b_0 \in B$.

\item \textbf{Generalized vs.\ strict inverses of Galois connections.} A monotone Galois connection $(f, g) : P \leftrightarrows Q$ need not be invertible as a map, nor is $(g, f)$ a monotone Galois connection from $Q$ to $P$ without passing to dual orders $(g, f) : Q^{\mathrm{op}} \leftrightarrows P^{\mathrm{op}}$. Instead, $g$ is the unique right adjoint (upper Galois inverse) of $f$, satisfying $f \circ g \circ f = f$ and $g \circ f \circ g = g$ and restricting to an order isomorphism $\mathrm{Fix}(g \circ f) \cong \mathrm{Fix}(f \circ g)$.
\end{enumerate}

\noindent\textbf{Salvaged sharp characterization.} On finite lattices, Clauses 1 and 2 hold canonically: every closure system $\mathcal{C} \subseteq L$ has a unique inclusion-least meet-generating set $\mathcal{M}(\mathcal{C})$ of length $|\mathcal{M}(\mathcal{C})|$ and a minimum concrete set-closure ground set of size $|\mathcal{J}(\mathcal{C})|$. Moreover, we prove the exact sharp dichotomy
\[
|\mathcal{M}(\mathcal{C})| = 0 \iff |\mathcal{C}| = 1, \qquad
|\mathcal{M}(\mathcal{C})| = 1 \iff |\mathcal{C}| = 2, \qquad
|\mathcal{M}(\mathcal{C})| \ge 2 \iff |\mathcal{C}| \ge 3,
\]
with equality $|\mathcal{M}(\mathcal{C})| = 2$ (and $|\mathcal{J}(\mathcal{C})| = 2$) attained up to isomorphism exclusively by the $3$-element chain $C_3$ and the $4$-element Boolean diamond $B_2$.

\section{Monotone Galois Connections and Galois Inverses}

\begin{definition}[Posets and Monotone Galois Connections]
Let $(P, \le_P)$ and $(Q, \le_Q)$ be partially ordered sets (posets). A \emph{monotone Galois connection} (or \emph{monotone adjoint pair}) between $P$ and $Q$, written $(f, g) : P \leftrightarrows Q$, consists of monotone maps $f : P \to Q$ (the \emph{left} or \emph{lower adjoint}) and $g : Q \to P$ (the \emph{right} or \emph{upper adjoint}, also called the \emph{upper Galois inverse} of $f$) such that
\begin{equation}\label{eq:adj}
\forall x \in P,\; \forall y \in Q,\qquad f(x) \le_Q y \iff x \le_P g(y).
\end{equation}
\end{definition}

\begin{proposition}[Fundamental Properties of Galois Inverses]\label{prop:galois-basic}
Let $(f, g) : P \leftrightarrows Q$ be a monotone Galois connection between posets.
\begin{enumerate}
\item \emph{(Unit--counit characterization)} A pair of maps $(f, g)$ is a monotone Galois connection if and only if $f$ and $g$ are monotone, $x \le_P g(f(x))$ for all $x \in P$ (inflation/unit), and $f(g(y)) \le_Q y$ for all $y \in Q$ (deflation/counit).
\item \emph{(Uniqueness of the Galois inverse)} If $(f, g_1)$ and $(f, g_2)$ are monotone Galois connections between $P$ and $Q$, then $g_1 = g_2$.
\item \emph{(Reflexive generalized inverse laws)} $f \circ g \circ f = f$ and $g \circ f \circ g = g$. Consequently, $c = g \circ f : P \to P$ is an idempotent closure operator ($x \le_P c(x) = c(c(x))$) with closed-element set
\[
\mathrm{Im}(g) = \{ g(y) \mid y \in Q \} = \{ x \in P \mid g(f(x)) = x \} = \mathrm{Fix}(g \circ f).
\]
\item \emph{(Order duality and restricted isomorphism)} $(g, f) : Q^{\mathrm{op}} \leftrightarrows P^{\mathrm{op}}$ is a monotone Galois connection between the dual posets, and $f|_{\mathrm{Fix}(g \circ f)}$ and $g|_{\mathrm{Fix}(f \circ g)}$ are mutually inverse order isomorphisms between $(\mathrm{Fix}(g \circ f), \le_P)$ and $(\mathrm{Fix}(f \circ g), \le_Q)$.
\end{enumerate}
\end{proposition}

\begin{proof}
(1) Setting $y = f(x)$ in~\eqref{eq:adj} yields $x \le_P g(f(x))$; setting $x = g(y)$ yields $f(g(y)) \le_Q y$. Conversely, if $f, g$ are monotone with $x \le_P g(f(x))$ and $f(g(y)) \le_Q y$, then $f(x) \le_Q y \implies x \le_P g(f(x)) \le_P g(y)$, and $x \le_P g(y) \implies f(x) \le_Q f(g(y)) \le_Q y$.

(2) For each $y \in Q$, $f(g_1(y)) \le_Q y$ gives $g_1(y) \le_P g_2(y)$ by~\eqref{eq:adj} for $(f, g_2)$, and symmetrically $g_2(y) \le_P g_1(y)$; antisymmetry gives $g_1(y) = g_2(y)$.

(3) Monotonicity of $f$ applied to $x \le_P g(f(x))$ gives $f(x) \le_Q f(g(f(x)))$, while deflation at $y = f(x)$ gives $f(g(f(x))) \le_Q f(x)$; hence $f(g(f(x))) = f(x)$. Dually, $g(f(g(y))) \le_P g(y)$ and $g(y) \le_P g(f(g(y)))$ yield $g(f(g(y))) = g(y)$. Thus $c(c(x)) = g(f(g(f(x)))) = g(f(x)) = c(x)$. If $x = g(y)$, then $g(f(x)) = g(f(g(y))) = g(y) = x$; conversely, if $g(f(x)) = x$, then $x = g(y)$ for $y = f(x)$.

(4) In $Q^{\mathrm{op}}$ and $P^{\mathrm{op}}$, $g(y) \le_{P^{\mathrm{op}}} x \iff x \le_P g(y) \iff f(x) \le_Q y \iff y \le_{Q^{\mathrm{op}}} f(x)$. For $x \in \mathrm{Fix}(g \circ f)$, $f(g(f(x))) = f(x)$ puts $f(x) \in \mathrm{Fix}(f \circ g)$ with $g(f(x)) = x$, and symmetrically for $y \in \mathrm{Fix}(f \circ g)$.
\end{proof}

\begin{remark}[Failure of Strict Invertibility and Un-dualized Symmetry]\label{rem:non-inv}
On the $2$-element chain $C_2 = \{0 < 1\}$, define $f_{\mathrm{const}}(x) = 0$ and $g_{\mathrm{const}}(y) = 1$. For all $x, y \in C_2$, $f_{\mathrm{const}}(x) = 0 \le y$ and $x \le 1 = g_{\mathrm{const}}(y)$ both hold unconditionally, so $(f_{\mathrm{const}}, g_{\mathrm{const}}) : C_2 \leftrightarrows C_2$ is a monotone Galois connection. However, $g_{\mathrm{const}}(0) = g_{\mathrm{const}}(1) = 1$ is not injective (no two-sided inverse exists), and the un-dualized swapped pair $(g_{\mathrm{const}}, f_{\mathrm{const}})$ is \emph{not} a monotone Galois connection on $C_2$ because $g_{\mathrm{const}}(1) = 1 \le 1$ holds whereas $1 \le 0 = f_{\mathrm{const}}(1)$ is false.
\end{remark}

\begin{example}[Non-Lattice Poset Obstruction $A_2$]\label{ex:antichain}
Let $A_2 = \{0, 1\}$ carry the discrete order $x \le y \iff x = y$. Then $(\mathrm{id}_{A_2}, \mathrm{id}_{A_2})$ is a monotone Galois connection on $A_2$ with $\mathrm{Im}(\mathrm{id}_{A_2}) = A_2$. However, $A_2$ has neither an empty meet (top element $\top$ satisfying $0 \le \top$ and $1 \le \top$) nor a binary meet $0 \wedge 1$. Thus $\mathrm{Im}(g)$ is not a meet-closed closure system on general posets unless ambient meets exist.
\end{example}

\begin{proposition}[Closure System Equivalence on Complete Lattices]\label{prop:complete-closure}
Let $L$ be a complete lattice (equipped with arbitrary infima $\bigwedge S$ for $S \subseteq L$, including $\top_L = \bigwedge \emptyset$). A subset $\mathcal{C} \subseteq L$ is a \emph{closure system} in $L$ ($\bigwedge S \in \mathcal{C}$ for every $S \subseteq \mathcal{C}$) if and only if $\mathcal{C} = \mathrm{Im}(g)$ for some monotone Galois connection $(f, g) : L \leftrightarrows Q$.
\end{proposition}

\begin{proof}
Given $(f, g) : L \leftrightarrows Q$ and $S \subseteq \mathrm{Im}(g) = \mathrm{Fix}(g \circ f)$, let $m = \bigwedge S$. For each $s \in S$, $m \le_L s$ implies $g(f(m)) \le_L g(f(s)) = s$, so $g(f(m)) \le_L \bigwedge S = m$. Together with $m \le_L g(f(m))$, antisymmetry yields $g(f(m)) = m \in \mathrm{Im}(g)$. Conversely, if $\mathcal{C} \subseteq L$ is a closure system, give $Q = \mathcal{C}$ the induced order, define $f(x) = \bigwedge \{ c \in \mathcal{C} \mid x \le_L c \} \in \mathcal{C}$, and let $g : \mathcal{C} \hookrightarrow L$ be the inclusion. Then $f(x) \le_L c \iff x \le_L c = g(c)$ for all $x \in L, c \in \mathcal{C}$, and $\mathrm{Im}(g) = \mathcal{C}$.
\end{proof}

\section{Taxonomy of Unit-Generated Closure Representations}

To ensure complete mathematical coverage, we formalize all four standard order-theoretic readings of a ``closure system generated by units'' and its ``representation length'' on a finite lattice $L$ with closure system $\mathcal{C} = \mathrm{Im}(g) \subseteq L$.

\subsection{Reading I: Minimal Meet-Generating Basis (Meet-Irreducible Units)}

For $B \subseteq L$, the closure system generated by $B$ under arbitrary meets in $L$ is $\langle B \rangle_{\bigwedge} = \{ \bigwedge S \mid S \subseteq B \}$ (with $\bigwedge \emptyset = \top_L$). An element $m \in \mathcal{C}$ is \emph{meet-irreducible in $\mathcal{C}$} if $m \neq \top_L$ and for all $a, b \in \mathcal{C}$, $a \wedge b = m \implies m = a \text{ or } m = b$. Denote by $\mathcal{M}(\mathcal{C})$ the set of meet-irreducible elements of $\mathcal{C}$.

\begin{theorem}[Unique Minimum Meet-Generating Representation]\label{thm:meet-irred}
Let $L$ be a finite lattice and $\mathcal{C} \subseteq L$ a closure system.
\begin{enumerate}
\item $\langle \mathcal{M}(\mathcal{C}) \rangle_{\bigwedge} = \mathcal{C}$.
\item For every subset $B \subseteq L$ satisfying $\langle B \rangle_{\bigwedge} = \mathcal{C}$, one has $\mathcal{M}(\mathcal{C}) \subseteq B \subseteq \mathcal{C}$.
\end{enumerate}
Consequently, $\mathcal{M}(\mathcal{C})$ is the unique inclusion-minimum (and unique minimum-cardinality) meet-generating set of $\mathcal{C}$, with representation length $\ell_{\mathrm{I}}(\mathcal{C}) = |\mathcal{M}(\mathcal{C})|$.
\end{theorem}

\begin{proof}
(1) Since $\mathcal{M}(\mathcal{C}) \subseteq \mathcal{C}$ and $\mathcal{C}$ is closed under arbitrary meets, $\langle \mathcal{M}(\mathcal{C}) \rangle_{\bigwedge} \subseteq \mathcal{C}$. For the reverse inclusion, we induct on $k(x) = |\{ y \in \mathcal{C} \mid x <_L y \}|$ for $x \in \mathcal{C}$. If $x = \top_L$, then $x = \bigwedge \emptyset \in \langle \mathcal{M}(\mathcal{C}) \rangle_{\bigwedge}$. If $x \neq \top_L$ and $x \in \mathcal{M}(\mathcal{C})$, then $x = \bigwedge \{x\} \in \langle \mathcal{M}(\mathcal{C}) \rangle_{\bigwedge}$. Otherwise $x = a \wedge b$ for some $a, b \in \mathcal{C}$ with $x <_L a$ and $x <_L b$. Then $\{ y \in \mathcal{C} \mid a <_L y \} \cup \{a\} \subseteq \{ y \in \mathcal{C} \mid x <_L y \}$, so $k(a) < k(x)$ and similarly $k(b) < k(x)$. By the induction hypothesis, $a = \bigwedge S_a$ and $b = \bigwedge S_b$ with $S_a, S_b \subseteq \mathcal{M}(\mathcal{C})$, whence $x = a \wedge b = \bigwedge (S_a \cup S_b) \in \langle \mathcal{M}(\mathcal{C}) \rangle_{\bigwedge}$.

(2) Suppose $\langle B \rangle_{\bigwedge} = \mathcal{C}$. For each $b \in B$, $b = \bigwedge \{b\} \in \mathcal{C}$, so $B \subseteq \mathcal{C}$. Let $m \in \mathcal{M}(\mathcal{C})$. Then $m = \bigwedge S$ for some finite subset $S = \{s_1, \dots, s_r\} \subseteq B \subseteq \mathcal{C}$. Since $m \neq \top_L$, $r \ge 1$. By induction on $r \ge 1$: if $r = 1$, $m = s_1 \in B$; if $r \ge 2$, $m = s_1 \wedge (s_2 \wedge \dots \wedge s_r)$ with both factors in $\mathcal{C}$, so meet-irreducibility forces $m = s_1 \in B$ or $m = s_2 \wedge \dots \wedge s_r$ (giving $m \in B$ by induction). Thus $\mathcal{M}(\mathcal{C}) \subseteq B$ and $|\mathcal{M}(\mathcal{C})| \le |B|$.
\end{proof}

\subsection{Reading II: Concrete Set-Closure Ground Set (Singleton Units)}

In Formal Concept Analysis and concrete closure theory, a closure system $\mathcal{C}$ is represented concretely as an intersection-closed family $\mathcal{F} \subseteq \mathcal{P}(U)$ over a finite ground set $U$ of singletons/units ($U \in \mathcal{F}$, and $A, B \in \mathcal{F} \implies A \cap B \in \mathcal{F}$), with $(\mathcal{C}, \le) \cong (\mathcal{F}, \subseteq)$. For a finite lattice $\mathcal{C}$, let $\mathcal{J}(\mathcal{C})$ denote its set of \emph{join-irreducible} elements ($j \neq \bot_{\mathcal{C}}$ and $a \vee_{\mathcal{C}} b = j \implies j = a \text{ or } j = b$). Notice that by Proposition~\ref{prop:galois-basic}(4), $\mathcal{J}(\mathcal{C})$ coincides under order duality with $\mathcal{M}_{Q^{\mathrm{op}}}(\mathrm{Im}(f))$!

\begin{theorem}[Minimum Concrete Set-Closure Ground Set]\label{thm:ground-set}
For every finite lattice $\mathcal{C}$, the minimum cardinality $|U|$ of a finite ground set $U$ admitting an intersection-closed family $\mathcal{F} \subseteq \mathcal{P}(U)$ with $(\mathcal{C}, \le) \cong (\mathcal{F}, \subseteq)$ is $\ell_{\mathrm{II}}(\mathcal{C}) = |\mathcal{J}(\mathcal{C})|$.
\end{theorem}

\begin{proof}
\emph{Upper bound:} Set $U = \mathcal{J}(\mathcal{C})$ and define $\varphi : \mathcal{C} \to \mathcal{P}(\mathcal{J}(\mathcal{C}))$ by $\varphi(x) = \{ j \in \mathcal{J}(\mathcal{C}) \mid j \le x \}$. Clearly $\varphi(\top_{\mathcal{C}}) = \mathcal{J}(\mathcal{C})$ and $\varphi(x \wedge y) = \varphi(x) \cap \varphi(y)$. Dualizing Theorem~\ref{thm:meet-irred}(1) shows every $x \in \mathcal{C}$ equals $\bigvee_{\mathcal{C}} \varphi(x)$, so $x \le y \iff \varphi(x) \subseteq \varphi(y)$. Hence $\mathcal{F} = \varphi(\mathcal{C})$ is a concrete closure system on $U$ isomorphic to $\mathcal{C}$ of size $|U| = |\mathcal{J}(\mathcal{C})|$.

\emph{Lower bound:} Let $\varphi : \mathcal{C} \to \mathcal{P}(U)$ be any order-embedding ($x \le y \iff \varphi(x) \subseteq \varphi(y)$) preserving binary meets ($\varphi(x \wedge y) = \varphi(x) \cap \varphi(y)$). For each $j \in \mathcal{J}(\mathcal{C})$, set $j_* = \bigvee_{\mathcal{C}} \{ x \in \mathcal{C} \mid x < j \}$. Clearly $j_* \le j$; if $j_* = j$, join-irreducibility and $j \neq \bot_{\mathcal{C}}$ would force $j = x$ for some $x < j$, a contradiction. Thus $j \not\le j_*$, so $\varphi(j) \not\subseteq \varphi(j_*)$, yielding a witness $u_j \in \varphi(j) \setminus \varphi(j_*) \subseteq U$. Suppose $u_{j_1} = u_{j_2} = u$ for $j_1, j_2 \in \mathcal{J}(\mathcal{C})$. Then $u \in \varphi(j_1) \cap \varphi(j_2) = \varphi(j_1 \wedge j_2)$. If $j_1 \wedge j_2 \neq j_1$, then $j_1 \wedge j_2 < j_1$ gives $j_1 \wedge j_2 \le (j_1)_*$, so $u \in \varphi((j_1)_*)$, impossible. Hence $j_1 \wedge j_2 = j_1$, and symmetrically $j_1 \wedge j_2 = j_2$, proving $j_1 = j_2$. Thus $j \mapsto u_j$ injects $\mathcal{J}(\mathcal{C})$ into $U$, giving $|\mathcal{J}(\mathcal{C})| \le |U|$.
\end{proof}

\subsection{Readings III and IV: Nonempty Meets, Closed Elements, and Unit Implications}

\begin{proposition}[Nonempty-Meet Generators and Closed-Element Count]\label{prop:reading-iii}
Let $L$ be a finite lattice and $\mathcal{C} \subseteq L$ a closure system.
\begin{enumerate}
\item If $\mathcal{C}$ is generated under \emph{nonempty} meets $\{ \bigwedge S \mid \emptyset \neq S \subseteq B \}$, then every generating set $B$ satisfies $\mathcal{M}(\mathcal{C}) \cup \{\top_L\} \subseteq B$, and $\mathcal{M}(\mathcal{C}) \cup \{\top_L\}$ generates $\mathcal{C}$. Thus the minimum nonempty-meet generator length is $\ell_{\mathrm{IIIa}}(\mathcal{C}) = |\mathcal{M}(\mathcal{C})| + 1$.
\item Counting closed elements directly (fixed points of the adjunction unit $x \mapsto g(f(x))$) gives $\ell_{\mathrm{IIIb}}(\mathcal{C}) = |\mathcal{C}| = |\mathrm{Im}(g)|$, which satisfies $|\mathcal{M}(\mathcal{C})| + 1 \le |\mathcal{C}| \le 2^{|\mathcal{M}(\mathcal{C})|}$.
\end{enumerate}
\end{proposition}

\begin{proof}
(1) If $\top_L = \bigwedge S$ with $S \neq \emptyset$, picking $s \in S$ gives $\top_L \le_L s \le_L \top_L$, so $\top_L = s \in B$. Moreover, any nonempty-meet generating set of $\mathcal{C}$ also generates $\mathcal{C}$ under all meets (since $\bigwedge \emptyset = \top_L \in \mathcal{C}$), so $\mathcal{M}(\mathcal{C}) \subseteq B$ by Theorem~\ref{thm:meet-irred}(2). Since $\top_L \notin \mathcal{M}(\mathcal{C})$, $|\mathcal{M}(\mathcal{C}) \cup \{\top_L\}| = |\mathcal{M}(\mathcal{C})| + 1 \le |B|$. Conversely, every $x \in \mathcal{C} \setminus \{\top_L\}$ is a nonempty meet from $\mathcal{M}(\mathcal{C})$, while $\top_L = \bigwedge \{\top_L\}$.

(2) Since $\top_L \in \mathcal{C} \setminus \mathcal{M}(\mathcal{C})$ and $\mathcal{M}(\mathcal{C}) \subseteq \mathcal{C}$, $|\mathcal{M}(\mathcal{C})| + 1 \le |\mathcal{C}|$; since $S \mapsto \bigwedge S$ surjects $\mathcal{P}(\mathcal{M}(\mathcal{C}))$ onto $\mathcal{C}$, $|\mathcal{C}| \le 2^{|\mathcal{M}(\mathcal{C})|}$.
\end{proof}

\begin{remark}[Reading IV: Unit-Implicational Bases]\label{rem:reading-iv}
In closure-system logic, a set-closure system $\mathcal{F} \subseteq \mathcal{P}(U)$ can also be specified as the models of a set $\Sigma$ of \emph{unit implications} $A \to \{b\}$ ($A \subseteq U$, $b \in U \setminus A$). Under this reading, for \emph{every} finite Boolean lattice $B_n = \mathcal{P}(\{1, \dots, n\})$ ($n \ge 0$) with $(\mathrm{id}_{B_n}, \mathrm{id}_{B_n})$, every subset is closed, so $\Sigma = \emptyset$ defines $B_n$ with minimum length $\ell_{\mathrm{IV}}(B_n) = 0 < 2$.
\end{remark}

\section{Counterexamples Refuting Conjecture 00000004507}

\subsection{Finite Lattice Counterexamples to the Lower Bound 2}

\begin{theorem}[Refutation of the Finite-Lattice Lower Bound $2$]\label{thm:finite-refutation}
The assertion that finite lattices attain a lower bound of $2$ on the closure-representation length of Galois inverses is false under every Reading~I--IV:
\begin{enumerate}
\item \emph{($1$-element finite lattice $C_1$)} Let $C_1 = \{0\}$ with $(\mathrm{id}_{C_1}, \mathrm{id}_{C_1})$. Then $C_1$ is a finite lattice with $\mathcal{C} = \{0\} = \{\top_{C_1}\}$, and
\[
\ell_{\mathrm{I}}(C_1) = 0, \quad \ell_{\mathrm{II}}(C_1) = 0, \quad
\ell_{\mathrm{IIIa}}(C_1) = \ell_{\mathrm{IIIb}}(C_1) = 1, \quad \ell_{\mathrm{IV}}(C_1) = 0,
\]
every one of which is strictly less than $2$.

\item \emph{(Nontrivial finite lattice $C_2$ with bijective identity Galois connection)} Let $C_2 = \{0 < 1\}$ ($|C_2| = 2$, $0 \neq 1$) with $(\mathrm{id}_{C_2}, \mathrm{id}_{C_2})$. Then $\mathcal{C} = C_2$ has $\mathcal{M}(C_2) = \{0\}$ and $\mathcal{J}(C_2) = \{1\}$ (concretely $C_2 \cong \{\emptyset, \{1\}\} \subseteq \mathcal{P}(\{1\})$), so
\[
\ell_{\mathrm{I}}(C_2) = 1 < 2, \qquad \ell_{\mathrm{II}}(C_2) = 1 < 2, \qquad \ell_{\mathrm{IV}}(C_2) = 0 < 2.
\]

\item \emph{(Constant Galois connections on arbitrary finite lattices)} For \emph{any} two finite lattices $L_1, L_2$ (including nontrivial lattices with $|L_1|, |L_2| \ge 2$), the maps $f(x) = \bot_{L_2}$ and $g(y) = \top_{L_1}$ form a monotone Galois connection $(f, g) : L_1 \leftrightarrows L_2$ with $\mathcal{C} = \mathrm{Im}(g) = \{\top_{L_1}\}$, for which $\ell_{\mathrm{I}}(\mathcal{C}) = \ell_{\mathrm{II}}(\mathcal{C}) = 0 < 2$ and $\ell_{\mathrm{IIIa}}(\mathcal{C}) = \ell_{\mathrm{IIIb}}(\mathcal{C}) = 1 < 2$.

\item \emph{(Non-constant Galois connections on arbitrarily large finite lattices)} Let $L$ be \emph{any} finite lattice with $|L| \ge 2$ and choose any $a \in L \setminus \{\top_L\}$. Define $f : L \to C_2$ and $g : C_2 \to L$ by
\[
f(x) = \begin{cases} 0 & \text{if } x \le_L a,\\ 1 & \text{if } x \not\le_L a, \end{cases}
\qquad
g(y) = \begin{cases} a & \text{if } y = 0,\\ \top_L & \text{if } y = 1. \end{cases}
\]
Then $(f, g) : L \leftrightarrows C_2$ is a monotone Galois connection with non-constant right adjoint ($g(0) = a \neq \top_L = g(1)$) and $\mathcal{C} = \mathrm{Im}(g) = \{a, \top_L\} \cong C_2$, whose unique minimal representation lengths under Readings~I and~II are $\ell_{\mathrm{I}}(\mathcal{C}) = \ell_{\mathrm{II}}(\mathcal{C}) = 1 < 2$.
\end{enumerate}
\end{theorem}

\begin{proof}
(1) In $C_1 = \{0\}$, $0 = \top_{C_1} = \bot_{C_1}$, so $\mathcal{M}(C_1) = \mathcal{J}(C_1) = \emptyset$ and $|C_1| = 1$.

(2) In $C_2 = \{0 < 1\}$, $1 = \top_{C_2}$ is excluded from $\mathcal{M}(C_2)$, while $0 \neq 1$ satisfies $a \wedge b = 0 \implies a = 0 \text{ or } b = 0$; hence $\mathcal{M}(C_2) = \{0\}$ has size $1$. Dually $\mathcal{J}(C_2) = \{1\}$ has size $1$.

(3) For all $x \in L_1$ and $y \in L_2$, $\bot_{L_2} \le_{L_2} y$ and $x \le_{L_1} \top_{L_1}$ both hold unconditionally, so $(f, g)$ is a monotone Galois connection with $\mathrm{Im}(g) = \{\top_{L_1}\} \cong C_1$.

(4) Clearly $f$ and $g$ are monotone. For $y = 0$, $f(x) \le 0 \iff f(x) = 0 \iff x \le_L a = g(0)$; for $y = 1$, $f(x) \le 1$ and $x \le_L \top_L = g(1)$ both hold for all $x \in L$. Thus $(f, g)$ is a monotone Galois connection with $\mathcal{C} = \mathrm{Im}(g) = \{a, \top_L\}$. Since $a \neq \top_L$, $\mathcal{M}(\mathcal{C}) = \{a\}$ and $\mathcal{J}(\mathcal{C}) = \{\top_L\}$, both of cardinality $1 < 2$.
\end{proof}

\begin{table}[ht]
\centering
\small
\setlength{\tabcolsep}{4.5pt}
\renewcommand{\arraystretch}{1.2}
\begin{tabular}{lccccc}
\hline
\textbf{Galois Connection Instance} & $|\mathcal{C}|$ & $\ell_{\mathrm{I}} = |\mathcal{M}(\mathcal{C})|$ & $\ell_{\mathrm{II}} = |\mathcal{J}(\mathcal{C})|$ & $\ell_{\mathrm{IIIa}} = |\mathcal{M}|+1$ & $\ell_{\mathrm{IIIb}} = |\mathcal{C}|$ \\
\hline
$C_1$, $(\mathrm{id}_{C_1}, \mathrm{id}_{C_1})$ & $1$ & $0$ & $0$ & $1$ & $1$ \\
Any $L_1 \leftrightarrows L_2$, $(f=\bot, g=\top)$ & $1$ & $0$ & $0$ & $1$ & $1$ \\
$C_2$, bijective $(\mathrm{id}_{C_2}, \mathrm{id}_{C_2})$ & $2$ & $1$ & $1$ & $2$ & $2$ \\
Any $|L| \ge 3$, $2$-element $L \leftrightarrows C_2$ & $2$ & $1$ & $1$ & $2$ & $2$ \\
$C_3$, $(\mathrm{id}_{C_3}, \mathrm{id}_{C_3})$ \emph{(witness)} & $3$ & $2$ & $2$ & $3$ & $3$ \\
$B_2$, $(\mathrm{id}_{B_2}, \mathrm{id}_{B_2})$ \emph{(witness)} & $4$ & $2$ & $2$ & $3$ & $4$ \\
\hline
\end{tabular}
\caption{Summary of closure-system representation lengths across key finite-lattice Galois connections.}
\label{tab:summary}
\end{table}

\subsection{Infinite Complete Lattice Counterexample to Minimal Generation}

We next show that Clause~1 of Conjecture~\texttt{00000004507} fails even for the identity Galois connection on an infinite complete distributive lattice. Let $\mathbb{D} = \{ (2m+1)/2^{k+1} \mid k \ge 0,\; 0 \le m < 2^k \} \subseteq (0, 1)$ denote the dyadic rationals in $(0, 1)$ with their strict linear order $<$, which is irreflexive, transitive, unbounded in $(0,1)$ ($\forall x \in \mathbb{D},\; \exists y_1, y_2 \in \mathbb{D},\; y_1 < x < y_2$), and dense ($\forall x < y \text{ in } \mathbb{D},\; \exists z \in \mathbb{D},\; x < z < y$).

Define $\mathrm{Cut}(\mathbb{D})$ to be the set of \emph{open rounded upper cuts} $U \subseteq \mathbb{D}$ satisfying:
\begin{enumerate}
\item \emph{(Upward closure)} $\forall x, y \in \mathbb{D},\; x \in U \land x < y \implies y \in U$;
\item \emph{(Openness / roundedness)} $\forall x \in U,\; \exists w \in U,\; w < x$.
\end{enumerate}
Order $\mathrm{Cut}(\mathbb{D})$ by \emph{reverse inclusion}: $U \le V \iff V \subseteq U$. Under this order, $\top = \emptyset$, $\bot = \mathbb{D}$, and for any family $\mathcal{S} \subseteq \mathrm{Cut}(\mathbb{D})$, the set union $\bigcup \mathcal{S}$ is upward-closed and rounded and forms the exact greatest lower bound $\bigwedge \mathcal{S}$ in $(\mathrm{Cut}(\mathbb{D}), \le)$! Arbitrary joins follow dually by $\bigvee \mathcal{S} = \bigwedge \{ W \mid \forall U \in \mathcal{S},\; U \le W \}$, making $\mathrm{Cut}(\mathbb{D})$ a complete lattice (order-isomorphic to $([0, 1], \le)$). For each $d \in \mathbb{D}$, the principal open cut $C_d = \{ x \in \mathbb{D} \mid d < x \}$ belongs to $\mathrm{Cut}(\mathbb{D})$ by transitivity and density.

\begin{theorem}[Complete Lattice Without Any Minimal Meet-Generating Set]\label{thm:infinite-cut}
Consider the monotone Galois connection $(\mathrm{id}, \mathrm{id})$ on the complete lattice $\mathrm{Cut}(\mathbb{D})$, whose closed-element closure system is $\mathcal{C} = \mathrm{Cut}(\mathbb{D})$.
\begin{enumerate}
\item \emph{(Non-vacuity)} The family $P = \{ C_d \mid d \in \mathbb{D} \}$ meet-generates $\mathrm{Cut}(\mathbb{D})$.
\item \emph{(No minimal meet-generating set)} For \emph{every} subset $B \subseteq \mathrm{Cut}(\mathbb{D})$ with $\langle B \rangle_{\bigwedge} = \mathrm{Cut}(\mathbb{D})$, there exists $b_0 \in B$ such that $\langle B \setminus \{b_0\} \rangle_{\bigwedge} = \mathrm{Cut}(\mathbb{D})$.
\item \emph{(Zero complete meet-irreducibles)} Every $U \in \mathrm{Cut}(\mathbb{D}) \setminus \{\top\}$ satisfies $U = \bigwedge S$ for a family $S \subseteq \mathrm{Cut}(\mathbb{D})$ with $U \notin S$.
\end{enumerate}
\end{theorem}

\begin{proof}
(1) For any $W \in \mathrm{Cut}(\mathbb{D})$, roundedness gives $x \in W \iff \exists d \in W,\; d < x \iff x \in \bigcup_{d \in W} C_d$, so $W = \bigwedge \{ C_d \mid d \in W \} \in \langle P \rangle_{\bigwedge}$.

(2) Let $\langle B \rangle_{\bigwedge} = \mathrm{Cut}(\mathbb{D})$. Unpacking $\bigwedge = \bigcup$, for every $U \in \mathrm{Cut}(\mathbb{D})$ and $x \in U$, there exists $V \in B$ with $U \le V$ (i.e.\ $V \subseteq U$) and $x \in V$. Since $3/4 \in C_{1/2}$, there exists some $b_0 \in B$ (taking $U = C_{1/2}$, $x = 3/4$). Put $B' = B \setminus \{b_0\}$. Given any $W \in \mathrm{Cut}(\mathbb{D})$ and $x \in W$, we claim there exists $V \in B'$ with $W \le V$ and $x \in V$ (whence $W = \bigwedge \{ V \in B' \mid W \le V \} \in \langle B' \rangle_{\bigwedge}$). By roundedness of $W$, pick $d \in W$ with $d < x$, so $W \le C_d$ and $x \in C_d$.
\begin{itemize}
\item \emph{Case 1:} $\neg (C_d \le b_0 \land x \in b_0)$. Applying $\langle B \rangle_{\bigwedge} = \mathrm{Cut}(\mathbb{D})$ to $C_d$ at $x \in C_d$ yields $V \in B$ with $C_d \le V$ and $x \in V$. Necessarily $V \neq b_0$, so $V \in B'$ and $W \le C_d \le V$.
\item \emph{Case 2:} $C_d \le b_0$ (meaning $b_0 \subseteq C_d$) and $x \in b_0$. By roundedness of $b_0$, pick $r \in b_0$ with $r < x$. Then $r \in b_0 \subseteq C_d$ gives $d < r$, so $C_d \le C_r$ and $x \in C_r$. Applying $\langle B \rangle_{\bigwedge} = \mathrm{Cut}(\mathbb{D})$ to $C_r$ at $x \in C_r$ yields $V \in B$ with $C_r \le V$ (i.e.\ $V \subseteq C_r$) and $x \in V$. Since $r \in b_0$ while $r \notin C_r$ (by irreflexivity $r \not< r$), $b_0 \not\subseteq C_r$, so $V \neq b_0$. Thus $V \in B'$ and $W \le C_d \le C_r \le V$.
\end{itemize}

(3) For any $U \neq \top$, $U = \bigwedge \{ C_d \mid d \in U \}$; for each $d \in U$, $d \in U \setminus C_d$, so $C_d \neq U$.
\end{proof}

\section{Salvaged Sharp Lower Bounds and Complete Classification}

We now prove the exact conditions under which the lower bound $2$ holds on finite lattices and classify all equality witnesses.

\begin{theorem}[Sharp Classification of Representation Lengths on Finite Lattices]\label{thm:classification}
Let $L$ be a finite lattice and let $\mathcal{C} = \mathrm{Im}(g) \subseteq L$ be the closed-element closure system of a monotone Galois connection $(f, g) : L \leftrightarrows Q$.
\begin{enumerate}
\item $|\mathcal{M}(\mathcal{C})| = 0 \iff |\mathcal{J}(\mathcal{C})| = 0 \iff |\mathcal{C}| = 1$.
\item $|\mathcal{M}(\mathcal{C})| = 1 \iff |\mathcal{J}(\mathcal{C})| = 1 \iff |\mathcal{C}| = 2$.
\item \emph{(Salvaged Sharp Lower Bound $2$)} $|\mathcal{M}(\mathcal{C})| \ge 2 \iff |\mathcal{J}(\mathcal{C})| \ge 2 \iff |\mathcal{C}| \ge 3$.
\item \emph{(Classification of Equality Witnesses)} One has $|\mathcal{M}(\mathcal{C})| = 2$ (equivalently $|\mathcal{J}(\mathcal{C})| = 2$) if and only if $(\mathcal{C}, \le_L)$ is order-isomorphic to either the $3$-element chain $C_3 = \{0 < 1 < 2\}$ or the $4$-element Boolean lattice $B_2 = \{0, 1\}^2$.
\end{enumerate}
\end{theorem}

\begin{proof}
By Proposition~\ref{prop:reading-iii}(2), $|\mathcal{M}(\mathcal{C})| + 1 \le |\mathcal{C}| \le 2^{|\mathcal{M}(\mathcal{C})|}$ always holds (and dually $|\mathcal{J}(\mathcal{C})| + 1 \le |\mathcal{C}| \le 2^{|\mathcal{J}(\mathcal{C})|}$ since $\mathcal{J}(\mathcal{C}) = \mathcal{M}(\mathcal{C}^{\mathrm{op}})$).
\begin{itemize}
\item If $|\mathcal{M}(\mathcal{C})| = 0$, then $1 \le |\mathcal{C}| \le 2^0 = 1$, so $|\mathcal{C}| = 1$; conversely, $|\mathcal{C}| = 1 \implies |\mathcal{M}(\mathcal{C})| + 1 \le 1 \implies |\mathcal{M}(\mathcal{C})| = 0$. This proves~(1).
\item If $|\mathcal{M}(\mathcal{C})| = 1$, then $2 \le |\mathcal{C}| \le 2^1 = 2$, so $|\mathcal{C}| = 2$; conversely, if $|\mathcal{C}| = 2$, then $|\mathcal{M}(\mathcal{C})| \le 1$ and $|\mathcal{M}(\mathcal{C})| \neq 0$ by~(1), so $|\mathcal{M}(\mathcal{C})| = 1$. Dualizing gives the same for $|\mathcal{J}(\mathcal{C})|$, proving~(2).
\item Since $|\mathcal{C}| \ge 1$ always, (1) and~(2) immediately imply $|\mathcal{M}(\mathcal{C})| \ge 2 \iff |\mathcal{C}| \ge 3 \iff |\mathcal{J}(\mathcal{C})| \ge 2$, proving~(3).
\item Suppose $|\mathcal{M}(\mathcal{C})| = 2$, say $\mathcal{M}(\mathcal{C}) = \{m_1, m_2\}$ with $m_1 \neq m_2$ and $m_1, m_2 <_L \top_L$. By Theorem~\ref{thm:meet-irred}(1), $\mathcal{C} = \{ \bigwedge S \mid S \subseteq \{m_1, m_2\} \} = \{\top_L, m_1, m_2, m_1 \wedge m_2\}$.
  \begin{itemize}
  \item If $m_1$ and $m_2$ are comparable, say $m_1 <_L m_2 <_L \top_L$, then $m_1 \wedge m_2 = m_1$, so $\mathcal{C} = \{m_1 <_L m_2 <_L \top_L\} \cong C_3$ (where $\mathcal{M}(C_3) = \{0, 1\}$ indeed has size $2$).
  \item If $m_1$ and $m_2$ are incomparable, then $m_1 \wedge m_2 <_L m_1, m_2 <_L \top_L$ are four distinct elements, so $\mathcal{C} = \{m_1 \wedge m_2, m_1, m_2, \top_L\} \cong B_2$ (where $\mathcal{M}(B_2) = \{(0,1), (1,0)\}$ indeed has size $2$).
  \end{itemize}
Since $C_3$ and $B_2$ are both self-dual, $|\mathcal{J}(\mathcal{C})| = |\mathcal{M}(\mathcal{C}^{\mathrm{op}})| = 2$ yields the exact same two isomorphism classes. \qedhere
\end{itemize}
\end{proof}

\section{Formal Verification in Core Lean 4 and Computational Audit}

All results in this paper have been machine-checked in self-contained Core Lean~4.19.0 (\nolinkurl{lean/Main.lean}, compiled with \nolinkurl{-DwarningAsError=true}, zero external dependencies, zero \nolinkurl{sorry}, zero \nolinkurl{native_decide}, and using only the standard foundational axioms \nolinkurl{propext}, \nolinkurl{Classical.choice}, and \nolinkurl{Quot.sound}):
{\raggedright
\begin{itemize}
\item Proposition~\ref{prop:galois-basic}: \nolinkurl{GaloisConnection.iff_unit_counit}, \nolinkurl{right_adjoint_unique}, \nolinkurl{f_g_f}, \nolinkurl{g_f_g}, \nolinkurl{mem_im_right_iff_closed}, \nolinkurl{dual}, and \nolinkurl{toClosed_toOpen}.
\item Example~\ref{ex:antichain} and Proposition~\ref{prop:complete-closure}: \nolinkurl{antichain2_no_top_and_no_meet}, \nolinkurl{galois_im_isClosureSystem}, and \nolinkurl{closureSystem_arises_from_galois}.
\item Theorem~\ref{thm:infinite-cut}: \nolinkurl{dyadicCut_no_minimal_generator}, \nolinkurl{dyadicCut_no_meetIrreducible}, and \nolinkurl{pCut_generates_cutLattice}.
\item Theorems~\ref{thm:meet-irred}--\ref{thm:ground-set} and Proposition~\ref{prop:reading-iii}: \nolinkurl{FiniteLattice.meetIrreds_generates}, \nolinkurl{meetIrreds_min_length}, \nolinkurl{concrete_ground_set_min_length}, and \nolinkurl{nonempty_meet_min_length}.
\item Theorem~\ref{thm:classification}(1)--(3): \nolinkurl{FiniteLattice.repLength_zero_iff}, \nolinkurl{repLength_one_iff}, and \nolinkurl{repLength_ge_two_iff}.
\item Theorem~\ref{thm:finite-refutation} and positive witnesses $C_3, B_2$: \nolinkurl{c1_refutes_lower_bound_two}, \nolinkurl{c2_bijective_refutes_lower_bound_two}, \nolinkurl{const_galois_refutes_all_readings}, \nolinkurl{nonconstant_two_elem_galois_length_one}, \nolinkurl{c3_and_b2_attain_lower_bound_two}, and master theorem \nolinkurl{conjecture_00000004507_disproved_and_classified}.
\end{itemize}
}

In addition, the standalone Python verification suite \texttt{verify.py} constructs all $78$ unlabeled finite lattices of cardinality $1 \le |L| \le 7$ (matching OEIS~A006966: $1, 1, 1, 2, 5, 15, 53$), exhaustively checks all $65$ monotone Galois connections across test posets, verifies Theorems~\ref{thm:meet-irred}, \ref{thm:ground-set}, and~\ref{thm:classification} by brute-force subset search over all $466$ closure systems for $|L| \le 6$, and confirms across all $78$ finite lattices up to $|L| = 7$ that $C_3$ and $B_2$ are the unique isomorphism classes with representation length $2$.

\begin{thebibliography}{9}
\bibitem{Birkhoff1967}
G.~Birkhoff, \emph{Lattice Theory}, 3rd ed., American Mathematical Society Colloquium Publications, Vol.~25, AMS, Providence, RI, 1967.

\bibitem{CaspardMonjardet2003}
N.~Caspard and B.~Monjardet, ``The lattices of closure systems, closure operators, and implicational systems on a finite set: a survey,'' \emph{Discrete Applied Mathematics} \textbf{127} (2003), no.~2, 241--269.

\bibitem{DaveyPriestley2002}
B.~A.~Davey and H.~A.~Priestley, \emph{Introduction to Lattices and Order}, 2nd ed., Cambridge University Press, Cambridge, 2002.

\bibitem{GanterWille1999}
B.~Ganter and R.~Wille, \emph{Formal Concept Analysis: Mathematical Foundations}, Springer-Verlag, Berlin, 1999.
\end{thebibliography}

\end{document}
